\section{USAMO 1987}
        \begin{enumerate}
        	\item (\textit{USAMO 1987}) Find all solutions to $(m^2+n)(m + n^2)= (m - n)^3$, where m and n are non-zero integers.


 



\textbf{Solution}

Expanding both sides, 
 \[m^3+mn+m^2n^2+n^3=m^3-3m^2n+3mn^2-n^3\]Note that $m^3$ can be canceled and as $n \neq 0$, $n$ can be factored out.Writing this as a quadratic equation in $n$: 
 \[2n^2+(m^2-3m)n+(3m^2+m)=0\].The discriminant $b^2-4ac$ equals 
 \[(m^2-3m)^2-8(3m^2+m)\]
 \[=m^4-6m^3-15m^2-8m\], which we want to be a perfect square.Miraculously, this factors as $m(m-8)(m+1)^2$. This is square \href{https://artofproblemsolving.com/wiki/index.php?title=Iff}{iff} $m^2-8m$ is square or $m+1=0$. It can be checked that the only nonzero $m$ that work are $-1, 8, 9$. Finally, plugging this in and discarding extraneous roots gives all possible ordered pairs $(m, n)$ as 
 \[\{(-1,-1),(8,-10),(9,-6),(9,-21)\}\].

 


	\item (\textit{USAMO 1987}) The feet of the angle bisectors of $\Delta ABC$ form a right-angled triangle. If the right-angle is at $X$, where $AX$ is the bisector of $\angle A$, find all possible values for $\angle A$.


 



\textbf{Solution}

Let two new segments $FD$ and $BE$ intersect at point $X$ and let $AD$ and $CF$ intersect at $Y$ .Consider quadrilateral $AFDC$. Then we get that the cross ratio $(B,Y;X,E)=-1$.Thus the condition $\angle FDE=90°$ implies $\angle FDA$=$\angle FDB$. Hence F is the A-excenter of the triangle $ADC$. Thus we get $\angle TAF$=$\angle FAD$=$\angle DAC$, where  $T$ is intersection of $FD$ and $AC$.Then we get $\angle BAC=120^\circ$

 


	\item (\textit{USAMO 1987}) $X$ is the smallest set of polynomials $p(x)$ such that: 


 
            \begin{enumerate}[label=(\roman*)]
                \item 1. $p = x$ belongs to $X$.
\item 2. If $r$ belongs to $X$, then $x\cdot r$ and $(x + (1 - x) \cdot r )$ both belong to $X$.
            \end{enumerate}
            Show that if $r(x)$ and $s(x)$ are distinct elements of $X$, then $r(x) \neq s(x)$ for any $0 < x < 1$. 


 



\textbf{Solution}

Let $s(x)$ be an arbitrary polynomial in $X.$ Then $0<s(x)<1$ when $0<x<1.$ Define $X_1=\{s(x)\in X:s(x)=x\cdot s_1(x)$ for some $s_1(x)\in X\},$ and $X_2=\{t(x)\in X: t(x)=x+(1-x)t_1(x)$ for some $t_1(x)\in X\}.$

 If $s(x) \in X_1$ and $t(x)\in X_2,$ we have $s(x) <x<t(x)$ for all $x$ with $0<x<1.$ Therefore $s(x)\ne t(x)$ for any $0<x<1.$

 For any $s(x), t(x) \in X_1$, Let $s(x)=x\cdot s_1(x)$ and $t(x)=x\cdot t_1(x)$ for $s_1(x), t_1(x) \in X.$ If $s_1(x) \ne t_1(x)$ for $0<x<1,$ then$s(x)-t(x)=x(s_1(x)-t_1(x))\ne 0$ for $0<x<1.$

 Similarly, for any $s(x), t(x) \in X_2$, Let $s(x)=x+(1-x) s_1(x)$ and $t(x)=x+(1-x) t_1(x)$ for $s_1(x), t_1(x) \in X.$ If $s_1(x) \ne t_1(x)$ for $0<x<1,$ then$s(x)-t(x)=(1-x)(s_1(x)-t_1(x))\ne 0$ for $0<x<1.$

 The proof is done by an induction.

 J.Z.

 


	\item (\textit{USAMO 1987}) Three circles $C_i$ are given in the plane: $C_1$ has diameter $AB$ of length $1$; $C_2$ is concentric and has diameter $k$ ($1<k<3$); $C_3$ has center $A$ and diameter $2k$. We regard $k$ as fixed. Now consider all straight line segments $XY$ which have one endpoint $X$ on $C_2$, one endpoint $Y$ on $C_3$, and contain the point $B$. For what ratio $\dfrac{XB}{BY}$ will the segment $XY$ have minimal length?


 



\textbf{Solution}

<i>This problem needs a solution. If you have a solution for it, please help us out by \href{https://artofproblemsolving.com/wiki/index.php?title=1987\_USAMO\_Problems/Problem\_4&action=edit}{adding it}.</i> 

 


	\item (\textit{USAMO 1987}) $a_1, a_2, \cdots, a_n$ is a sequence of 0's and 1's. T is the number of triples $(a_i, a_j, a_k)$ with $i<j<k$ which are not equal to (0, 1, 0) or (1, 0, 1). For $1\le i\le n$, $f(i)$ is the number of $j<i$ with $a_j = a_i$ plus the number of $j>i$ with $a_j\neq a_i$. Show that $T=\sum_{i=1}^n f(i)\cdot\left(\frac{f(i)-1}2\right)$. If n is odd, what is the smallest value of T?


 



\textbf{Solution}

Let $g(i)$ be the number of $j<i$ with $a_j = a_i$ and $h(i)$ be the number of $j>i$ with $a_j\neq a_i$, so that $f(i) = g(i) + h(i)$. Then
 \[f(i)\left(\frac{f(i)-1}2\right) = (g(i)+h(i))\left(\frac{g(i)+h(i)-1}2\right)\]
 \[\frac{g(i)^2+h(i)^2+2g(i)h(i)-g(i)-h(i)}{2}\]
 \[\frac{g(i)(g(i)-1)}{2} + \frac{h(i)(h(i)-1)}{2} + g(i)h(i)\]
 \[\binom{g(i)}{2} + \binom{h(i)}{2} + \binom{g(i)}{1}\binom{h(i)}{1}\]

 The first term counts the number of triples of the form $(0,0,0)$ or $(1,1,1)$ with $a_i$ as the last member. The second term counts the number of triples of the form $(0,1,1)$ or $(1,0,0)$ with $a_i$ as the first member. The third term counts the number of triples of the form $(0,0,1)$ or $(1,1,0)$ with $a_i$ as the middle member. Thus $\sum_{i=1}^n f(i)\cdot\left(\frac{f(i)-1}2\right)$ exhaustively and uniquely counts all triples except for those of the forms $(1,0,1)$ and $(0,1,0)$.

 
Finding the smallest value of T is trickier. First, consider what happens when we reverse the order of the sequence. For a given $a_i$, the less than/greater than relation between $a_i$ and $a_j$, which determines whether $f$ counts $a_i = a_j$ or $a_i \neq a_j$, is reversed for all $i \neq j$, so in the new sequence $a_i$ will correspond to a function value of $n-1-f(i)$. However, notice that T does not change when we reverse the order of the sequence, as any triple $(1,0,1)$ or $(0,1,0)$ will also be present in the reversed sequence. We'll set T from the original sequence equal to T from the reversed sequence

 
 \[\frac{1}{2}\sum_{i=1}^n f(i)(f(i)-1)) = \frac{1}{2}\sum_{i=1}^n (n-1-f(i))\cdot\left((n-2-f(i))\right)\]
 \[\sum_{i=1}^n(f(i)^2 - f(i)) = \sum_{i=1}^n (n^2 - 2n -nf(i)-n+2+f(i)-nf(i)+2f(i)+f(i)^2)\]
 \[-\sum_{i=1}^n f(i) = \sum_{i=1}^n ((3-2n)f(i)+n^2-3n+2)\]
 \[-\sum_{i=1}^n f(i) = (3-2n)\sum_{i=1}^n (f(i))+n^3-3n^2+2n\]
 \[-\sum_{i=1}^n f(i) = (3-2n)\sum_{i=1}^n (f(i))+n(n-1)(n-2)\]
 \[(2n-4)\sum_{i=1}^n f(i) = n(n-1)(n-2)\]
 \[\sum_{i=1}^n f(i) = \frac{n(n-1)}{2}\]

 We find that for all sequences, $\sum_{i=1}^n f(i) = \frac{n(n-1)}{2}$. We'll use the method of Lagrange multipliers to minimize $\frac{1}{2}\sum_{i=1}^n f(i)(f(i)-1)$ subject to $\sum_{i=1}^n f(i) = \frac{n(n-1)}{2}$. Note that this only establishes a sufficient condition for a minimum, we will show later that it is possible to define a sequence which satisfies that condition.

 
 \[\mathcal{L} (f(1),...,f(n), \lambda) = \frac{1}{2}\sum_{i=1}^n f(i)(f(i)-1) + \lambda\left(\sum_{i=1}^n(f(i))-\frac{n(n-1)}{2}\right)\]

 For all $i$, we get $\frac{d\mathcal{L}}{df(i)} = \frac{1}{2}(2f(i)-1) + \lambda$. Setting the gradient to 0 and solving quickly yields $f(i) = \frac{1}{2}-\lambda$. Then we take the derivative with respect to $\lambda$: $\frac{d\mathcal{L}}{d\lambda} = \sum_{i=1}^n (f(i)) - \frac{n(n-1)}{2}$. Setting the gradient to 0 and substituting our expression for $f(i)$ gives $0 = n(\frac{1}{2}-\lambda) - \frac{n(n-1)}{2}$, which simplifies to $\lambda = \frac{2-n}{2}$. Finally, substituting this back into the expression for $f(i)$ gives $f(i) = \frac{n-1}{2}$. 

 
This tells us that if we can define a sequence for any odd $n$ such that $f(i) = \frac{n-1}{2}$ for all $i$, this would certainly be a sequence which minimizes $T$. In fact, an alternating sequence $0,1,0,1,...,0$ satisfies this, because for any $i$, the number of $j<i$ such that $a_j = a_i$ is $\lfloor \frac{i-1}{2} \rfloor$ and the number of $j>i$  such that $a_j \neq a_i$ is $\lceil \frac{n-i}{2} \rceil$. For odd $i$, $\lfloor \frac{i-1}{2} \rfloor + \lceil \frac{n-i}{2} \rceil = \frac{i-1}{2} + \frac{n-i}{2} = \frac{n-1}{2}$ and for even $i$, $\lfloor \frac{i-1}{2} \rfloor + \lceil \frac{n-i}{2} \rceil = \frac{i-2}{2} + \frac{n-i+1}{2} = \frac{n-1}{2}$. Thus T is minimized by this alternating sequence, where it is equal to $\sum_{i=1}^n \frac{n-1}{2} \frac{n-3}{4} = \frac{n(n-1)(n-3)}{8}$.

 


 


\end{enumerate}